% !TeX spellcheck = en_US

\section{Appendix B: Optical and Instrumental Anomaly Detection Module} \label{sec:ap:anom}

This appendix describes the methodology for detecting instrumental anomalies that are undocumented in the official metadata of the \textit{Euclid} and \textit{DESI} missions. The approach is based on the hypothesis that data corruptions (such as chromatic filter misalignment, poorly subtracted background gradients, hot pixels, or fiber saturation) represent extreme deviations from natural physical and statistical laws in the manifold of the multimodal latent space.

\subsection{Mahalanobis Distance in Latent Space}

To quantify how anomalous the joint latent representation of an object is, it is assumed that the distribution of normal embeddings $\mathbf{e}_i$ can be locally approximated by a multivariate Gaussian distribution. 

We define the mean vector of the pretraining sample $\boldsymbol{\mu} \in \mathbb{R}^d$ and the covariance matrix $\boldsymbol{\Sigma} \in \mathbb{R}^{d \times d}$ as:

\begin{equation}
\boldsymbol{\mu} = \frac{1}{M}\sum_{i=1}^M \mathbf{e}_i
\label{eq:ap_mean}
\end{equation}

\begin{equation}
\boldsymbol{\Sigma} = \frac{1}{M-1}\sum_{i=1}^M (\mathbf{e}_i - \boldsymbol{\mu})(\mathbf{e}_i - \boldsymbol{\mu})^T
\label{eq:ap_cov}
\end{equation}

For a new object with joint embedding $\mathbf{x} \in \mathbb{R}^d$, its anomaly score is defined by calculating the Mahalanobis distance $D_M(\mathbf{x})$:

\begin{equation}
D_M(\mathbf{x}) = \sqrt{(\mathbf{x} - \boldsymbol{\mu})^T \boldsymbol{\Sigma}^{-1} (\mathbf{x} - \boldsymbol{\mu})}
\label{eq:ap_mahalanobis}
\end{equation}

Since the latent dimension $d = 768$ is high and the covariance matrix can be singular or unstable due to the collinearity of latent features, a Ledoit-Wolf shrinkage regularization was applied to the covariance matrix:

\begin{equation}
\boldsymbol{\Sigma}_{\text{reg}} = (1 - \alpha) \boldsymbol{\Sigma} + \alpha \frac{\text{Tr}(\boldsymbol{\Sigma})}{d} \mathbf{I}
\label{eq:ap_regularized_cov}
\end{equation}

where $\alpha \in [0, 1]$ is the optimal shrinkage coefficient and $\mathbf{I}$ is the identity matrix.

Objects with $D_M(\mathbf{x}) > \tau$ (where the threshold $\tau$ is defined according to the $99.5\%$ percentile of the training sample) are classified as anomalous candidates for manual review or purging.

\subsection{Typology of Identified Anomalies}

Visual analysis of the sources classified with the highest anomaly scores revealed three major families of instrumental systematics not described in the official photometric or spectroscopic control tables:

\begin{enumerate}
	\item \textbf{Chromatic and registration misalignment in Euclid:} 
	Objects presenting a physical spatial offset of several pixels in the cutout centering between the VIS camera and the near-infrared NISP channels. This generates a false artificial color gradient (for example, a red elliptical galaxy laterally offset from its blue halo in VIS), a physically impossible scenario for normal isolated galaxies.
	
	\item \textbf{Saturation and fiber contamination in DESI:}
	Spectra where a defective pixel in the spectrograph detector produces extremely narrow artificial emission lines (one pixel Doppler width) with unphysical fluxes that exceed the continuum by orders of magnitude. The model detects that these transitions do not respect natural quantum atomic flux ratios (for example, the constant theoretical atomic flux ratio $F([\text{O\,III}]\,\lambda 5007) / F([\text{O\,III}]\,\lambda 4959) \approx 3$).
	
	\item \textbf{Spurious background gradients:}
	Euclid cutouts located near bright Milky Way stars that suffer from scattered light or internal optical reflections (\textit{ghosts}), introducing brightness ramps or chromatic contamination that biases automated non-parametric morphological estimates.
\end{enumerate}

This module demonstrates that multimodal self-supervised pretraining implicitly assimilates normal instrumental correlations, allowing the latent space to be used as an unsupervised quality auditor that is far superior to rigid logical rules on catalogs.


\begin{figure}[h]
	\centering
	\begin{subfigure}[b]{\textwidth}
		\centering
		\includegraphics[width=\textwidth]{example-image-a}
		\caption{}
		\label{fi:ap:art}
	\end{subfigure}
	\hfill
	\begin{subfigure}[b]{\textwidth}
		\centering
		\includegraphics[width=\textwidth]{example-image-b}
		\caption{}
		\label{fi:ap:ano}
	\end{subfigure}
	\caption[]{}
	\label{fi:ap:artefacts}
\end{figure}
